\documentclass[11pt]{article}

\usepackage[margin=1.1in]{geometry}
\usepackage{amsmath,amssymb,amsthm}
\usepackage{mathtools}
\usepackage{stmaryrd}
\usepackage{microtype}
\usepackage[T1]{fontenc}
\usepackage{palatino}
\usepackage{enumitem}
\usepackage{array}
\usepackage[hidelinks]{hyperref}

\theoremstyle{definition}
\newtheorem{definition}{Definition}
\newtheorem{principle}{Principle}
\theoremstyle{plain}
\newtheorem{proposition}{Proposition}
\newtheorem{conjecture}{Conjecture}
\newtheorem{question}{Question}
\theoremstyle{remark}
\newtheorem{remark}{Remark}

\newcommand{\inn}[3]{\mathrm{for}(#1 \leftarrow #2)\,#3}
\newcommand{\nil}{\mathbf{0}}
\newcommand{\rcong}{\equiv}
\newcommand{\rbisim}{\sim_{r}}
\newcommand{\sem}[1]{\llbracket #1 \rrbracket}
\newcommand{\AC}{\mathrm{AC}}
\newcommand{\DC}{\mathrm{DC}}
\newcommand{\CC}{\mathrm{CC}}

\title{\bfseries From the Axiom of Choice to Fair Scheduling\\[3pt]
\large Toward a Curry--Howard Correspondence for Nondeterminism,\\
in which choice principles are the types and fair schedulers are the terms}

\author{L. Gregory Meredith\\
\small F1R3FLY.io / F1R3FLY Industries\\
\small \texttt{lgreg.meredith@gmail.com}}

\date{Research Note --- June 2026}

\begin{document}
\maketitle

\begin{abstract}
A companion note \cite{consensus} proposed that Gabbay's Coalition Logic stands to the
rho calculus as linear logic stands to the $\pi$-calculus: a declarative logic whose
proofs, under a Curry--Howard term assignment, extract correct-by-construction process
implementations of consensus protocols, with cut elimination as message passing. This
note proposes the same shape one level deeper, for nondeterminism itself. We argue that
\emph{the mathematical theory of choice} --- the axiom of choice and its gradations,
ordered by the Weihrauch lattice --- is a declarative language (a logic, a type theory)
for nondeterminism, and that \emph{the theory of fairness and scheduling} is the
corresponding term language: choice principles are the types, schedulers and fairness
disciplines are the terms, and running a schedule is normalisation. The bridge already
exists in fragments. The constructive axiom of choice is a theorem whose proof
\emph{is} a choice function (Martin-L\"of); classical realizability interprets dependent
choice and full choice by programs equipped with control and ``quote'' primitives
(Krivine); bar recursion and the selection monad give choice an operational, monadic
content (Spector; Berardi--Bezem--Coquand; Escard\'o--Oliva). What is missing, and what
we outline, is the assembly of these into a correspondence for \emph{concurrent}
nondeterminism, whose realisers are schedulers. The opening is precise: the
proofs-as-processes tradition works by design in the confluent, deadlock-free fragment
(Wadler's CP is Church--Rosser), excluding exactly the races the mobile process calculi
exist to model; choice-logic is the complementary logic of that discarded
nondeterminism. Under the proposed correspondence the strength of a choice principle is
the altitude of its scheduler in the hypercomputation fibration of \cite{agency}, so
that \emph{subject reduction becomes the statement that scheduling cannot bootstrap a
stronger oracle} --- and the central open hinge of \cite{agency}, whether interference
can crack open a determinate system, becomes precisely the question whether the term
assignment satisfies subject reduction. We close with an extraction pipeline parallel to
the consensus one, and with its composition: a protocol correct-by-construction in its
logic and fairly-scheduled-by-construction in its dynamics, both powers typed.
\end{abstract}

\section{Introduction}

Three artefacts now sit side by side. A reflective higher-order (rho) calculus \cite{Mer05}
with a keyed-store (RSpace) semantics in which parallel composition is the cut and
communication is the ``forgiving zip'' of dual bags at a channel \cite{qcs,paths}. An essay
\cite{agency} arguing that the nondeterminism of that zip is not missing information but a
\emph{resource} --- a choice, in the exact mathematical sense, graded in strength by the
axiom of choice and climbing, with that strength, the stalks of a fibration of
hypercomputation over the category of computational models. And a research note \cite{consensus}
proposing that Coalition Logic proofs extract rho implementations of consensus protocols by a
Curry--Howard term assignment, with cut elimination as the protocol's message passing.

The note \cite{consensus} instantiates a general pattern: a \emph{declarative logic}, an
\emph{operational term language}, and an \emph{extraction} that turns proofs into programs by
normalisation. The present note observes that the resource of \cite{agency} --- choice --- is
itself an instance of that pattern, and develops it. The declarative side is the mathematical
theory of choice: a choice principle is a specification of what nondeterministic resolution is
available or required. The operational side is the theory of fairness and scheduling: a
scheduler is a program that resolves nondeterminism, a fairness discipline a constraint on
schedulers. The claim, stated as a slogan and then unpacked, is:
\begin{quote}
\emph{Choice is the logic of scheduling. A choice principle is the type of a scheduler; a
scheduler is a realiser of a choice principle; running the schedule is normalisation.}
\end{quote}
Why this matters, beyond symmetry. Schedulers and fairness disciplines are today written by
hand and verified, if at all, after the fact. A Curry--Howard correspondence would make them
\emph{correct-by-construction}: extracted from a specification of the nondeterminism a system
requires. More pointedly, it would make their computational power \emph{typed}. The thesis of
\cite{agency} is that a scheduler can smuggle in an oracle --- that the resolution of
nondeterminism can import computational strength a system did not otherwise have. A type
discipline in which the choice principle is the type is exactly an accounting of how much
strength a given scheduler smuggles in: the type is the altitude, and a well-typed scheduler
is one whose imported power is declared.

We proceed as the companion note does. Section~\ref{sec:two} lays out the two theories.
Section~\ref{sec:bridge} records the bridge that already exists, in the realizability and
selection-function literature. Section~\ref{sec:gap} identifies the gap --- the confluence of
proofs-as-processes --- which is the opening this programme fills. Section~\ref{sec:corr}
sketches the correspondence and its term assignment. Section~\ref{sec:pipeline} gives the
extraction pipeline and its composition with the consensus pipeline. Sections~\ref{sec:conn}
and \ref{sec:open} place the work and list what a rigorous development must discharge.

\section{Two theories of nondeterminism}\label{sec:two}

\subsection{Choice, declaratively}

The mathematical theory of choice is a graded family of principles asserting that a section
exists for a family of nonempty sets. We read each principle as a \emph{type}: a specification
of the nondeterministic resolution a computation is permitted to assume.

\begin{itemize}[leftmargin=1.4em]
\item \emph{Finite choice.} For a finite family the section exists with no axiom, constructively;
and when the family is \emph{presented} --- each fibre given as a list, as a finite bag of
prefixes is --- the choice function is definable outright. This is the type of a scheduler that
needs no special power: the term is its own choice structure.
\item \emph{Countable choice $\CC$ and dependent choice $\DC$.} A countable family, or a family
whose later fibres depend on earlier selections. $\DC$ is the type of a scheduler with memory and
lookahead, choosing each step in the context of the residuals of the last.
\item \emph{Full $\AC$.} An arbitrary family. The type of a scheduler that can survey an
unindexed contention and select coherently --- the strongest principle, and (Diaconescu
\cite{diaconescu}) the one that, in a topos, imports excluded middle.
\end{itemize}

\noindent
These principles are not linearly ordered by logical strength alone; the finer structure is the
\emph{Weihrauch lattice} \cite{weihrauch}, whose objects are choice (and choice-like) principles
ordered by uniform reducibility and composed by the operation ``solve one, then, using its answer,
the next.'' We take the Weihrauch lattice to be the entailment order of the logic of choice: the
order in which one scheduler-type is at least as strong as another, and the composition along which
schedules are chained. That it is a \emph{lattice} and not a chain is the declarative shadow of the
``jumble, not tower'' ontology of \cite{agency}.

\begin{remark}[Constructive versus classical choice is a polarity, not a quarrel]
In intuitionistic type theory the choice principle for a presented family is a theorem, because
a proof of $\forall x\,\exists y\,\varphi$ already \emph{is} a function $x \mapsto y$ (Martin-L\"of
\cite{martinlof}); its realiser is a pure term. The set-theoretic, non-constructive $\AC$ is a
different type, realised only by programs with classical control. The two are not rival accounts of
one principle but two principles of different strength --- exactly the grading above, with the
constructive one at the base of the stalk and the classical one above.
\end{remark}

\subsection{Fairness and scheduling, operationally}

A scheduler resolves the nondeterminism a system leaves open: at each point of contention it
selects which enabled interaction fires. A \emph{fairness discipline} constrains the schedulers
under consideration --- weak fairness (a continuously enabled action eventually fires), strong
fairness (an infinitely-often-enabled action eventually fires), and their refinements --- and is,
classically, the assumption under which liveness and progress can be guaranteed (Francez
\cite{francez}). We take this to be the term language: schedulers are the terms, fairness
disciplines the typing constraints they must satisfy. The polarities of the scheduling literature
--- \emph{demonic} (adversarial: must survive every resolution) and \emph{angelic} (cooperative:
may seek a good resolution) --- are, in this reading, the two ways a term can inhabit a choice-type,
and (as \cite{agency} conjectures) the left and right adjoints to reindexing along the strength
order.

\begin{remark}[The three-valued zone is where fairness lives]
The companion note's three truth values --- running ($t$), blocked/faulty ($b$), dead ($f$)
\cite{consensus} --- have a scheduling reading. A choice point is \emph{live} (a redex, schedulable),
\emph{blocked} (enabled but not yet fired, a pending consume), or \emph{dead} (no admissible
resolution). Fairness is precisely the discipline of the blocked zone: the guarantee that a choice
point which stays live does not stay blocked forever. So three-valuedness is not an ornament; it is
where the term language does its characteristic work, and a fairness type is a constraint on the
$b$-zone.
\end{remark}

\section{The bridge already exists, in fragments}\label{sec:bridge}

The correspondence is not built from nothing. The computational content of choice has been studied
for decades, and in every case the realiser of a choice principle is a program whose extra power
matches the principle's strength --- which is exactly the typing discipline we want, in a sequential
setting.

\begin{itemize}[leftmargin=1.4em]
\item \emph{Constructive $\AC$ as a theorem (Martin-L\"of \cite{martinlof}).} The proof is the
choice function; the term assignment is the identity insight. This is the base case of the
correspondence: at the bottom of the stalk, the scheduler is read straight off the proof.
\item \emph{Bar recursion (Spector \cite{spector}).} Spector's bar recursion realises the
double-negation shift and hence countable choice and classical analysis, as a recursion schema in an
extension of G\"odel's System~T. A choice principle is realised by a \emph{specific recursion operator}
--- a term with a definite computational shape, a scheduler with backward-inductive lookahead.
\item \emph{Computational content of $\AC$ (Berardi--Bezem--Coquand \cite{bbc}).} A realiser of the
axiom of choice that proceeds by \emph{learning} --- updating a finite approximation to the choice
function as it goes. Read operationally it is a scheduler that commits provisionally and revises:
precisely the on-demand, never-completed choice process \cite{agency} calls for at $\omega$.
\item \emph{Dependent choice by control (Krivine \cite{krivine}).} Krivine's classical realizability
interprets $\DC$ using a ``quote'' instruction and a clock --- the realiser is a program with a
specific non-functional primitive. The strength of the principle is the strength of the primitive:
the cleanest existing statement that \emph{a choice-type's inhabitant is a program with matching
extra power}.
\item \emph{The selection monad (Escard\'o--Oliva \cite{escardo}).} Selection functions ---
``given a predicate, choose a witnessing point'' --- form a monad whose algebra is the algebra of
choice-as-program, with bar recursion as its characteristic operation and backward induction as its
paradigm. This is the categorical home for ``scheduler as term'': the selection monad is, we propose,
the monad of the scheduler calculus.
\end{itemize}

\noindent
Each of these gives choice a term content. None gives it in a \emph{concurrent} setting, where the
contention is a race at a channel and the realiser is a scheduler over a store. Assembling them there
is the programme.

\section{The gap: proofs-as-processes is confluent}\label{sec:gap}

The companion note's analogy was ``Coalition Logic is to rho as linear logic is to $\pi$.'' The second
half of that analogy --- the proofs-as-processes programme of Abramsky and of Bellin--Scott
\cite{abramsky,bellinscott}, and the propositions-as-sessions of Caires--Pfenning and Wadler
\cite{cairespfenning,wadler} --- is exactly where the present opening lies. In that programme a
linear-logic proof is a session-typed process and cut elimination is communication; it is a genuine
Curry--Howard correspondence for concurrency, and it is, by design, \emph{confluent and deadlock-free}.
Wadler's CP is Church--Rosser: cut elimination has a unique normal form, sessions do not race, and
deadlock is ruled out by the logic.

That confluence is the point and the limitation. It means the proofs-as-processes correspondence captures
the \emph{communication skeleton} --- who talks to whom on which protocol --- while discarding precisely
the nondeterminism the mobile process calculi exist to express: the races of \cite{agency}, two sends
contending for one receive, the choice of which meets which. A confluent system cannot type that choice,
because confluence is the assertion that the choice does not matter. So linear logic types the structure
of interaction and is silent on its resolution.

\begin{principle}[Choice-logic is the complement of linear logic]\label{prin:complement}
The Curry--Howard content of concurrent computation factors into two orthogonal axes. Linear logic
(propositions-as-sessions) types the \emph{confluent communication skeleton}: the deterministic,
deadlock-free protocol structure. The theory of choice types the \emph{nondeterministic resolution}:
which of the admissible interactions, at a genuine race, actually fires. The first has a developed
Curry--Howard correspondence; the second is the one this note proposes. Their term languages are,
respectively, session-typed processes and fair schedulers, and a full account of a mobile process is
the pairing of the two.
\end{principle}

\noindent
This is why the theory of choice, and not linear logic, is the right declarative partner for the
nondeterministic content. Linear logic deliberately excluded it; choice is the logic of exactly what
was excluded.

\section{The proposed correspondence}\label{sec:corr}

We sketch the term assignment. As in \cite{consensus} this is a programme, stated to be discharged, not
a completed system.

\subsection{The dictionary}

\begin{center}
\renewcommand{\arraystretch}{1.25}
\begin{tabular}{>{\raggedright}p{0.46\textwidth} >{\raggedright\arraybackslash}p{0.46\textwidth}}
\textbf{Theory of choice (declarative)} & \textbf{Fairness and scheduling (operational)}\\[2pt]
\hline
choice principle / formula & scheduler discipline / type\\
a proof of $\forall x\,\exists y\,\varphi$ & a choice/selection function (the term)\\
the family $\prod_k X_k$ of admissible pairings & the schedule (a section of the family)\\
finite, presented choice & pure scheduler (no special primitive)\\
countable / dependent choice $\DC$ & bar recursion; scheduler with memory/lookahead\\
full $\AC$ & oracle scheduler; strongest control\\
strength in the Weihrauch lattice & altitude in the hypercomputation stalk \cite{agency}\\
cut elimination / normalisation & running the schedule (resolving the race)\\
subject reduction & scheduling preserves altitude (no oracle bootstrap)\\
progress & a live choice is schedulable (no spurious deadlock)\\
three-valued enabledness ($t/b/f$) & fairness governs the blocked ($b$) zone\\
Diaconescu: $\AC \Rightarrow$ excluded middle & strongest scheduler imports classical control\\
\hline
\end{tabular}
\end{center}

\subsection{Formulae, proofs, and the site of the term}

At an aperture --- a channel carrying both polarities, a race --- the obligation is a $\forall\exists$:
for the contended slot, \emph{there exists} an admissible pairing. The formula is the choice-type of that
aperture; a proof is a scheduler that, for the slot, produces the pairing; the term acts at exactly the
forgiving zip of \cite{qcs}, where the underlying lists are chosen and matched. The denotation $\sem{-}$
stays a function and the zip is the separate dynamics \cite{qcs}; in the present reading the zip's choice
of list orderings is the \emph{term} inhabiting the aperture's choice-type, and normalisation --- running
the schedule --- is the firing of matched pairs with surplus remaining.

A confluent interior --- a macro-component of \cite{agency} --- is, in this light, a region whose choice-type
is \emph{proof-irrelevant}: several pairings inhabit it but all normalise to $\rbisim$-equal results, so the
term carries no observable information. The genuine proof content sits at the boundary apertures. ``Capacity
lives at the apertures'' \cite{agency} becomes ``the term content lives where the choice-type is not
proof-irrelevant.''

\subsection{Subject reduction is the no-bootstrap law --- and the open hinge}

The two metatheorems one demands of a term assignment have, here, direct operational meaning.

\emph{Progress}: a well-typed live choice can be scheduled --- the system does not deadlock except where the
type (the $f$-zone) says no resolution exists. This is the fairness guarantee, internalised as a typing property.

\emph{Subject reduction}: running the schedule preserves the type --- the scheduler does not, by reducing,
come to realise a \emph{stronger} choice principle than it was typed for. Operationally: \emph{scheduling
cannot bootstrap an oracle}. This is the proof-theoretic form of the no-implicit-interaction discipline of
\cite{paths}, and of the constant-altitude claim of \cite{agency}: a scheduler of degree $\mathbf d$ keeps
degree $\mathbf d$ as it runs.

\begin{question}[The hinge, as subject reduction]\label{q:sr}
Does the term assignment satisfy subject reduction across the cascade in which one interaction manufactures a
deeper one (the In-deep step of \cite{paths})? Equivalently --- and this is the open hinge of \cite{agency} in
proof-theoretic dress --- can the resolution of a boundary race cause the interior to acquire a choice-type it
did not have, climbing the stalk as it reduces?
\end{question}

\noindent
If subject reduction holds, the factorisation of \cite{agency} is stable: determinate interiors stay determinate
under scheduling, and a scheduler's imported power is fixed by its type once and for all. If it fails, the
failure \emph{is} perturbation: a high-capacity scheduler, reducing against a low-capacity system, manufactures
new choice in it --- a well-typed term reducing to an ill-typed one, the type-theoretic signature of an agent
changing another's capacity. The same lemma settles the agency question and the metatheory of the correspondence.

\section{The synthesis pipeline}\label{sec:pipeline}

Mirroring \cite{consensus}, the extraction has three stages.

\emph{Stage 1 --- specification.} The nondeterminism a system requires is specified as a choice-type: which
contentions must be resolvable, under which fairness discipline, at what permitted strength. The permitted
strength is a Weihrauch-degree bound --- a ceiling on the oracle the scheduler may import.

\emph{Stage 2 --- proof and extraction.} A proof that the specification is satisfiable in the logic of choice is
constructed; the term assignment extracts a scheduler. It is correct-by-construction: it resolves exactly the
specified contentions, under the specified fairness, \emph{and within the specified strength} --- the type
certifies that it smuggles in no more oracle than declared.

\emph{Stage 3 --- optimisation.} As in \cite{consensus}, extraction yields a naive scheduler; bisimulation-preserving
transformations reduce overhead while preserving the realised choice-type, and the phlogiston metering of the
rholang runtime serves as the objective, now with the added invariant that optimisation must not raise the
scheduler's altitude.

\begin{remark}[The two pipelines compose]
Consensus protocols need fair scheduling to make progress; the companion note extracts the protocol's
communication (the \emph{what}), this note extracts the fair resolution of its nondeterminism (the \emph{how}).
Stacked, they yield a protocol correct-by-construction in its Coalition-Logic specification \emph{and}
fairly-scheduled-by-construction in its dynamics, with both the protocol logic and the scheduling power typed ---
and, by Question~\ref{q:sr}, with a guarantee (if subject reduction holds) that running the protocol does not
covertly raise the computational strength its scheduler was licensed.
\end{remark}

\section{Connections and context}\label{sec:conn}

\emph{To the agency essay \cite{agency}.} This note is its proof-theoretic face. The aperture is the site of a
non-trivial choice-type; the macro-component is a proof-irrelevant region; the hypercomputation stalk is the
strength order made into types; interference is a realiser supplied from a stronger type than the lower system's
own; and the open hinge is subject reduction (Question~\ref{q:sr}).

\emph{To the consensus note \cite{consensus}.} Both instantiate one pattern --- declarative logic, term language,
extraction by normalisation. The notes compose rather than merely resemble: Coalition Logic over the communication
structure, choice-logic over the nondeterministic resolution, with the rho calculus the common term substrate and
cut-elimination-as-communication the shared engine.

\emph{To proofs-as-processes \cite{abramsky,bellinscott,cairespfenning,wadler}.} The complementary axis
(Principle~\ref{prin:complement}). Linear logic types the confluent skeleton; choice-logic types the discarded
nondeterminism. A mobile process is fully accounted for only by both, and the present programme is the missing half.

\emph{To realizability and selection functions \cite{martinlof,spector,bbc,krivine,escardo}.} The technical
grounding: choice principles already have realisers whose power tracks their strength. The contribution sought here
is to carry that into a concurrent, store-based setting where the realisers are schedulers and the contention is a
race --- with the selection monad \cite{escardo} the likely categorical backbone.

\section{Open problems}\label{sec:open}

\begin{enumerate}[label=P\arabic*.,leftmargin=2.4em]
\item \textbf{The logic of choice as a formal system.} Present the choice principles as a deductive calculus with
the Weihrauch lattice \cite{weihrauch} as its entailment order and Weihrauch composition $\star$ as the cut. Identify
the analogue of geometric/structural rules for the gradations (finite, $\CC$, $\DC$, $\AC$).
\item \textbf{The scheduler calculus.} Define the term language precisely --- schedulers as terms over the RSpace
store, fairness disciplines as typing constraints, the selection monad \cite{escardo} as the computational structure
--- and the operation that runs a schedule.
\item \textbf{Term assignment and metatheory.} Define the assignment of Section~\ref{sec:corr} and prove progress
(fairness as a typing property) and subject reduction (Question~\ref{q:sr}).
\item \textbf{The grading theorem.} Prove that the choice-type of an extracted scheduler equals its
Weihrauch/Turing degree equals its altitude in the fibration of \cite{agency}: type, degree, and altitude are one.
\item \textbf{The hinge.} Resolve Question~\ref{q:sr} --- whether scheduling can climb the stalk --- proving stability
of the factorisation if subject reduction holds, or characterising perturbation as its definite failure mode if it does
not.
\item \textbf{Composition with consensus.} Make precise the stacking of this pipeline with \cite{consensus}: a single
extraction whose output is correct-by-construction in protocol logic and fair-and-bounded-by-construction in scheduling.
\item \textbf{Mechanisation.} Formalise the logic of choice, the scheduler calculus, and the term assignment in a proof
assistant (Lean~4), enabling machine-checked extraction of bounded-strength fair schedulers.
\end{enumerate}

\section{Conclusion}

The companion note \cite{consensus} read Coalition Logic as a declarative language whose proofs extract rho
implementations of consensus protocols. This note reads the mathematical theory of choice as a declarative language
for nondeterminism, whose proofs --- under a term assignment still to be built --- extract fair schedulers. The two
notes are instances of one pattern and compose into one pipeline; together they say that both the \emph{what} of a
concurrent protocol and the \emph{how} of resolving its nondeterminism can be specified declaratively and synthesised
correct-by-construction. The realizability tradition has long shown that choice principles have realisers whose power
tracks their strength; the proofs-as-processes tradition has long shown that confluent concurrency is Curry--Howard.
What lies between --- the Curry--Howard content of \emph{non}-confluent concurrency, the logic and term language of the
race --- is the theory of choice paired with the theory of scheduling, and its central metatheorem, subject reduction,
is the same lemma that decides whether agency is a fixed allotment or something agents make in one another. If the
correspondence holds, a scheduler is a proof, its fairness is a type, and the oracle it imports is written on its face.

\begin{thebibliography}{99}

\bibitem[Mer05]{Mer05} L. G. Meredith and M. Radestock. A reflective higher-order calculus.
\emph{Electr.\ Notes Theor.\ Comput.\ Sci.}, 141(5):49--67, 2005.

\bibitem[Agcy]{agency} L. G. Meredith. \emph{Where Does Agency Live? Nondeterminism as a computational resource,
and agency across the computational and biological axes}. Manuscript, 2026.

\bibitem[Cons]{consensus} L. G. Meredith. \emph{From Coalition Logic to Rho Calculus: Toward a Curry--Howard
Correspondence for Consensus Protocol Synthesis}. Research Note, F1R3FLY, 2026.

\bibitem[QCS]{qcs} L. G. Meredith. \emph{Quoting is Colour-Swap: a model of the rho calculus in the knotted
universe}. Manuscript, 2026.

\bibitem[Path]{paths} L. G. Meredith. \emph{Paths are Subspaces: a cut-triggered polymorphism of RSpace over
path-keys, and its distributive law}. Manuscript, 2026.

\bibitem[ML84]{martinlof} P. Martin-L\"of. \emph{Intuitionistic Type Theory}. Bibliopolis, Naples, 1984.

\bibitem[Spe62]{spector} C. Spector. Provably recursive functionals of analysis: a consistency proof of analysis by
an extension of principles formulated in current intuitionistic mathematics. In \emph{Recursive Function Theory},
Proc.\ Symp.\ Pure Math.\ V, pages 1--27. AMS, 1962.

\bibitem[BBC98]{bbc} S. Berardi, M. Bezem, and T. Coquand. On the computational content of the axiom of choice.
\emph{J. Symbolic Logic}, 63(2):600--622, 1998.

\bibitem[Krv03]{krivine} J.-L. Krivine. Dependent choice, `quote' and the clock. \emph{Theoret.\ Comput.\ Sci.},
308(1--3):259--276, 2003.

\bibitem[EO10]{escardo} M. H. Escard\'o and P. Oliva. Selection functions, bar recursion and backward induction.
\emph{Math.\ Structures Comput.\ Sci.}, 20(2):127--168, 2010.

\bibitem[Fra86]{francez} N. Francez. \emph{Fairness}. Texts and Monographs in Computer Science. Springer, 1986.

\bibitem[Abr94]{abramsky} S. Abramsky. Proofs as processes. \emph{Theoret.\ Comput.\ Sci.}, 135(1):5--9, 1994.

\bibitem[BS94]{bellinscott} G. Bellin and P. J. Scott. On the $\pi$-calculus and linear logic. \emph{Theoret.\
Comput.\ Sci.}, 135(1):11--65, 1994.

\bibitem[CP10]{cairespfenning} L. Caires and F. Pfenning. Session types as intuitionistic linear propositions. In
\emph{CONCUR 2010}, LNCS 6269, pages 222--236. Springer, 2010.

\bibitem[Wad14]{wadler} P. Wadler. Propositions as sessions. \emph{J. Funct.\ Programming}, 24(2--3):384--418, 2014.

\bibitem[BGP21]{weihrauch} V. Brattka, G. Gherardi, and A. Pauly. Weihrauch complexity in computable analysis. In
\emph{Handbook of Computability and Complexity in Analysis}, pages 367--417. Springer, 2021.

\bibitem[Dia75]{diaconescu} R. Diaconescu. Axiom of choice and complementation. \emph{Proc.\ Amer.\ Math.\ Soc.},
51:176--178, 1975.

\end{thebibliography}

\end{document}
